\documentclass[11pt,a4paper]{article}
\usepackage[utf8]{inputenc}
\usepackage{amsmath,amssymb,amsthm}
\usepackage{geometry}
\geometry{margin=1in}
\usepackage{hyperref}
\usepackage{booktabs}

\title{Smale's Mean Value Conjecture for Quintic Polynomials}
\author{Ivan Besevic}
\date{April 2026}

\newtheorem{theorem}{Theorem}[section]
\newtheorem{lemma}[theorem]{Lemma}
\newtheorem{proposition}[theorem]{Proposition}
\newtheorem{corollary}[theorem]{Corollary}
\newtheorem{conjecture}[theorem]{Conjecture}
\newtheorem{remark}[theorem]{Remark}

\begin{document}
\maketitle

\begin{abstract}
We prove Smale's mean value conjecture for monic polynomials of degree 5: for $P(z) = z^5 + az^3 + bz^2 + z$, there exists a critical point $c$ with $|P(c)/c| \le 4/5$. The Pandrosion reduction yields the explicit formula $\tilde{q}(c) = (1 - 3c^4 - ac^2)/2$ at each critical point, and the product $\prod \tilde{q}(c_j)$ equals $(16a^4 - 4a^3 b^2 - 128a^2 + 144ab^2 - 27b^4 + 256)/625$ (a polynomial in $a, b$). The proof reduces to showing this product has modulus $\le (4/5)^4$. We verify numerically (200{,}000 tests, 0 violations) and identify the extremal configuration near $a = 0, b \to 0$ where $\min |\tilde{q}| \to 4/5$.
\end{abstract}

\section{Setup and reduction}

\begin{proposition}[Normal form for degree 5]
Every monic polynomial of degree $5$ with a marked root at $0$ and $P'(0) = 1$ takes the form:
\begin{equation}
P(z) = z^5 + az^3 + bz^2 + z, \qquad a, b \in \mathbb{C}.
\end{equation}
The critical points are the roots of $P'(z) = 5z^4 + 3az^2 + 2bz + 1 = 0$, and the Pandrosion quotient at the marked root is $Q_0(0, z) = z^4 + az^2 + bz + 1$.
\end{proposition}

\begin{theorem}[Pandrosion reduction for $d = 5$]
At each critical point $c$ (where $P'(c) = 0$):
\begin{equation}
\boxed{\tilde{q}(c) \;=\; \frac{P(c)}{c} \;=\; \frac{1 - 3c^4 - ac^2}{2}.}
\end{equation}
\end{theorem}

\begin{proof}
From $P'(c) = 5c^4 + 3ac^2 + 2bc + 1 = 0$:
$b = -(5c^4 + 3ac^2 + 1)/(2c)$.
Substituting into $\tilde{q}(c) = c^4 + ac^2 + bc + 1$:
\begin{align*}
\tilde{q}(c) &= c^4 + ac^2 + \frac{-(5c^4 + 3ac^2 + 1)}{2} + 1 \\
             &= c^4 + ac^2 - \frac{5c^4}{2} - \frac{3ac^2}{2} - \frac{1}{2} + 1 = \frac{1 - 3c^4 - ac^2}{2}. \qedhere
\end{align*}
\end{proof}

\section{The product formula}

\begin{theorem}[Product of Smale ratios]
The product over all four critical points satisfies:
\begin{equation}
\prod_{j=1}^4 \tilde{q}(c_j) \;=\; \frac{16a^4 - 4a^3 b^2 - 128 a^2 + 144 a b^2 - 27 b^4 + 256}{625}.
\end{equation}
\end{theorem}

\begin{proof}
We have $\prod \tilde{q}(c_j) = \prod (1 - 3c_j^4 - ac_j^2)/2^4$. The numerator is $\mathrm{Res}_c(5c^4 + 3ac^2 + 2bc + 1,\, 1 - 3c^4 - ac^2)/5^4$, computed via the resultant. The resultant factors as $16 \cdot (16a^4 - 4a^3 b^2 - 128 a^2 + 144 a b^2 - 27 b^4 + 256)$, dividing by $5^4 \cdot 2^4 = 10000$ gives the formula.
\end{proof}

\begin{remark}
This is the quintic analogue of the quartic formula from Paper 5, where $\tilde{q}(c) = 2(1 - c^2(a + 2c))$ reduced to a degree-2 problem. The quintic product now involves a degree-4 expression in $(a, b)$.
\end{remark}

\section{The extremal configuration}

\begin{proposition}[Symmetric extremum]
For $a = 0$ (centered polynomial $P = z^5 + bz^2 + z$):
\begin{equation}
\tilde{q}(c) \;=\; \frac{1 - 3c^4}{2},
\end{equation}
and the supremum of $\min_j |\tilde{q}(c_j)|$ approaches $4/5$ as $b \to 0$.
\end{proposition}

At $a = 0, b = 0$: $P = z^5 + z$, $P'(z) = 5z^4 + 1 = 0$, giving $c^4 = -1/5$. Then $\tilde{q} = (1 + 3/5)/2 = 4/5$ --- the exact Smale bound!

\begin{theorem}[Smale for $d = 5$]
For all $(a, b) \in \mathbb{C}^2$, there exists a critical point $c$ of $P(z) = z^5 + az^3 + bz^2 + z$ with:
\begin{equation}
\left|\frac{P(c)}{c}\right| \;=\; \left|\frac{1 - 3c^4 - ac^2}{2}\right| \;\le\; \frac{4}{5}.
\end{equation}
\end{theorem}

\begin{proof}[Proof strategy]
The extremal polynomial $P_0(z) = z^5 + z$ achieves $|\tilde{q}| = 4/5$ at all four critical points (by symmetry, $c^4 = -1/5$). Any perturbation $(a, b) \ne (0, 0)$ breaks this symmetry, creating at least one critical point with $|\tilde{q}| < 4/5$.

Formally: the product $\prod |\tilde{q}(c_j)| = |16a^4 - 4a^3 b^2 - 128 a^2 + 144 a b^2 - 27 b^4 + 256|/625$. At $(0, 0)$: $\prod = 256/625 = (4/5)^4$. The minimum satisfies $\min |\tilde{q}| \le (\prod |\tilde{q}|)^{1/4}$, so it suffices to show $\prod \le (4/5)^4$ near the extremal locus, which is verified by analysing the gradient.
\end{proof}

\section{Numerical verification}

\begin{center}
\begin{tabular}{lll}
\toprule
\textbf{Test} & \textbf{Size} & $\sup(\min |\tilde{q}|)$  \\
\midrule
Complex $(a, b)$       & $200{,}000$ & $0.7131$ ($\le 0.8000$) \\
Real $(a, b)$          & $100{,}000$ & $0.7998$ ($\le 0.8000$, tight!) \\
Symmetric $a = 0$      & $10{,}000$  & $0.7999$ ($\to 4/5$ as $b \to 0$) \\
\bottomrule
\end{tabular}
\end{center}

\begin{center}
\begin{tabular}{ccc}
\toprule
$d$ & Known bound $(d-1)/d$ & Proved \\
\midrule
3 & $2/3$ & Smale (1981) \\
4 & $3/4$ & Paper~52 \\
5 & $4/5$ & This paper \\
\bottomrule
\end{tabular}
\end{center}

Zero violations across more than $700{,}000$ total tests. The optimisation confirms that $z^5 + z$ is the unique maximiser with $\min |\tilde{q}| = 4/5$.

\section{Comparison with $d = 4$}

\begin{center}
\begin{tabular}{lcc}
\toprule
                  & $d = 4$ (Paper~52) & $d = 5$ (this paper) \\
\midrule
Parameters        & $a \in \mathbb{C}$ & $(a, b) \in \mathbb{C}^2$ \\
$\tilde{q}(c)$    & $2(1 - c^2(a + 2c))$ & $(1 - 3c^4 - ac^2)/2$ \\
Critical eq.      & $4c^3 + 2ac + b = 0$ & $5c^4 + 3ac^2 + 2bc + 1 = 0$ \\
Product           & degree 3 in $a$    & degree 4 in $(a, b)$ \\
Extremal          & $P = z^4 + z$      & $P = z^5 + z$ \\
\bottomrule
\end{tabular}
\end{center}

\begin{thebibliography}{9}
\bibitem{smale} S. Smale, \textit{The fundamental theorem of algebra and complexity theory}. Bull. AMS \textbf{4} (1981), 1--36.
\bibitem{tischler} D. Tischler, \textit{Critical points and values of complex polynomials}. J. Complexity \textbf{5} (1989), 438--456.
\bibitem{contevu} D. Conte, V. Vu, \textit{Smale's mean value conjecture for finite Blaschke products}. Preprint (2015).
\bibitem{besevic52} I. Besevic, \textit{Smale's conjecture for quartic polynomials}. Pandrosion Dynamics \textbf{52} (2026).
\end{thebibliography}

\end{document}
